\thispagestyle{empty}
\begin{center}
\vspace*{2.2cm}
{\color{poprule}\rule{\linewidth}{1.2pt}}\\[10pt]
{\sffamily\bfseries\color{popink}\Huge Poplog and Pop-11}\\[12pt]
{\sffamily\Large\color{popaccent} A Working Introduction}\\[10pt]
{\color{poprule}\rule{\linewidth}{1.2pt}}\\[24pt]
{\sffamily\large\color{popgrey}
 The two-level virtual machine, the core language,\\
 the four hosted front-ends, and the C interface}\\[40pt]
\end{center}

\vfill
\noindent
{\sffamily\small\color{popgrey}
This book accompanies the Poplog source tree at
\texttt{github.com/IoTone/poplog} --- a maintained fork of the Poplog system
originally developed at the University of Sussex.  Every Pop-11, Forth,
Prolog, Common Lisp and Standard ML listing in it was executed against that
tree on Apple Silicon (macOS, \texttt{arm64}) while the text was written;
the outputs shown are the outputs obtained.\par
\vspace{8pt}
Full in-tree documentation --- 900+ \textsc{help}, \textsc{teach} and
\textsc{ref} files --- is browsable at
\texttt{https://iotone.github.io/poplog/}.\par
\vspace{14pt}
\textcopyright{} The Poplog contributors.  Poplog is free software; this
document is distributed with the source tree under the same terms.\par
}
\vspace{10pt}
\clearpage
